\section{המשפט האופטי וגלים חלקיים}

\begin{summary}{מבוא: חשיבות ניסויי פיזור}
ניסויי פיזור הם כלי מרכזי בחקר מבנה החומר, החל מניסוי רתרפורד שגילה את הגרעין ועד למאיצי חלקיקים מודרניים (כמו ב-\LR{LHC}). הרעיון הוא להאיץ חלקיקים, לירות אותם על מטרה (או חלקיק אחר), ומתוך תבנית הפיזור ללמוד על הכוחות והמבנה הפנימי.
\end{summary}

\subsection{המשפט האופטי (\LR{The Optical Theorem})}

המשפט האופטי הוא תוצאה מפתיעה המקשרת בין גודל הנמדד ב"קדימה" (\LR{Forward Scattering}) לבין חתך הפעולה הכולל. מקור המשפט באופטיקה (לורד ריילי, המאה ה-19) והוא נובע משימור הסתברות.

\begin{theorembox}{המשפט האופטי}
חתך הפעולה הכולל ($\sigma_{tot}$) פרופורציונלי לחלק המדומה של אמפליטודת הפיזור בזווית אפס ($f(0)$):
\[
\sigma_{tot} = \frac{4\pi}{k} \text{Im}[f(0)]
\]
כאשר $f(0)$ היא אמפליטודת הפיזור בכיוון ההתקדמות המקורי ($\theta=0$).
\end{theorembox}


נוכיח את המשפט באמצעות אופרטור ה-\LR{$T$-matrix} ומשוואת ליפמן-שווינגר.

\textbf{שלב 1: הגדרות} \\
אמפליטודת הפיזור מוגדרת על ידי:
\[
f(\theta) = -\frac{m}{2\pi \hbar^2} \langle k' | T | k \rangle
\]
כאשר עבור פיזור קדימה $k=k'$. אופרטור $T$ מקיים $T|k\rangle = V|\psi_{in}\rangle$, כאשר $|\psi_{in}\rangle$ הוא המצב המלא (הכולל את הפיזור).

\textbf{שלב 2: שימוש במשוואת ליפמן-שווינגר} \\
אנו מעוניינים לחשב את $\text{Im}\langle k|T|k \rangle$. נשתמש בקשר האופרטורי:
\[
T|k\rangle = V|\psi_{in}\rangle
\]
נציב את משוואת ליפמן-שווינגר (בצורתה ההפוכה עבור הברה והקט):
\[
\text{Im}\langle k|T|k \rangle = \text{Im}\langle \psi_{in} | V | \psi_{in} \rangle + \text{Im} \langle \psi_{in} | V \frac{1}{E - H_0 + i\epsilon} V | \psi_{in} \rangle
\]
האיבר הראשון $\langle \psi_{in} | V | \psi_{in} \rangle$ הוא ממשי (כי $V$ הרמיטי), ולכן החלק המדומה שלו מתאפס.

\textbf{שלב 3: שימוש בזהות מתמטית} \\
נשתמש בזהות עבור פונקציית הגרין (\LR{Cauchy Principal Value}):
\[
\frac{1}{x + i\epsilon} = \mathcal{P}\left(\frac{1}{x}\right) - i\pi \delta(x)
\]
החלק של ה-$\mathcal{P}$ (הערך העיקרי) הוא ממשי ולכן נופל בחישוב ה-$\text{Im}$. נשארנו עם:
\[
\text{Im}\langle k|T|k \rangle = -\pi \langle \psi_{in} | V \delta(E - H_0) V | \psi_{in} \rangle
\]

\textbf{שלב 4: הכנסת בסיס שלם} \\
נכניס בסיס שלם של מצבי תנע חופשיים $\int d^3k' |k'\rangle\langle k'|$:
\[
\text{Im}\langle k|T|k \rangle = -\pi \int d^3k' |\langle k'|V|\psi_{in}\rangle|^2 \delta\left(E - \frac{\hbar^2 k'^2}{2m}\right)
\]
נזכור ש-$\langle k'|V|\psi_{in}\rangle = \langle k'|T|k\rangle$. נבצע אינטגרציה על הגודל $k'$ בעזרת הדלתא, ונישאר עם אינטגרל על הזוויות:
\[
\text{Im}\langle k|T|k \rangle \propto -\int d\Omega |f(\theta)|^2
\]
האינטגרל על הזוויות הוא בדיוק ההגדרה של $\sigma_{tot}$, וכך מתקבל המשפט האופטי (עד כדי הקבועים).


\subsection{פיתוח לגלים חלקיים (\LR{Partial Wave Analysis})}

במקום לטפל בגל המישורי ובגל המתפזר כגוש אחד, אנו מפרקים אותם לרכיבים של תנע זוויתי ($l$). הסיבה היא שבפוטנציאל מרכזי, התנע הזוויתי נשמר.

\begin{formulabox}{פיתוח גל מישורי לגלים כדוריים (\LR{Rayleigh Expansion})}
גל מישורי בכיוון $z$ ניתן לכתיבה כסכום של גלים כדוריים:
\[
e^{ikz} = e^{ikr\cos\theta} = \sum_{l=0}^{\infty} i^l (2l+1) j_l(kr) P_l(\cos\theta)
\]
כאשר:
\begin{itemize}
    \item $j_l(kr)$ הן פונקציות בסל ספריות (רגולריות בראשית).
    \item $P_l(\cos\theta)$ הם פולינומי לז'נדר (בגלל הסימטריה הצירית, $m=0$).
\end{itemize}
\end{formulabox}

\begin{notebox}{התנהגות אסימפטוטית}
במרחקים גדולים ($r \to \infty$), פונקציית בסל מתנהגת כמו סינוס:
\[
j_l(kr) \approx \frac{\sin(kr - l\pi/2)}{kr}
\]
זה מייצג סופרפוזיציה של גל כדורי נכנס וגל כדורי יוצא \textbf{ללא הפרש פאזה} (באין פוטנציאל).
\end{notebox}

\subsubsection{השפעת הפוטנציאל והיסט הפאזה }
כאשר קיים פוטנציאל (קצר טווח), הגלים הכדוריים הנכנסים לא משתנים (הם מגיעים מהאינסוף), אך הגלים היוצאים עוברים שינוי. בפיזור אלסטי, השינוי היחיד האפשרי הוא בפאזה.

\begin{definitionbox}{מטריצת הפיזור $S_l$ והיסט הפאזה}
אנו מגדירים את השינוי בגל היוצא בעזרת המקדם $S_l$:
\[
\psi_{total} \xrightarrow{r\to\infty} \text{\LR{Incoming}} + S_l \cdot \text{\LR{Outgoing}}
\]
עבור פיזור אלסטי (שימור הסתברות בכל גל חלקי):
\[
S_l = e^{2i\delta_l}
\]
כאשר $\delta_l$ הוא \textbf{היסט הפאזה} (\LR{Phase Shift}).
\end{definitionbox}

\begin{resultbox}{אמפליטודת הפיזור באמצעות גלים חלקיים}
על ידי חיסור הגל המישורי המקורי מפונקציית הגל המלאה, מקבלים את פונקציית הגל המתפזרת, וממנה את $f(\theta)$:
\[
f(\theta) = \frac{1}{2ik} \sum_{l=0}^{\infty} (2l+1) (S_l - 1) P_l(\cos\theta)
\]
או באמצעות היסט הפאזה:
\[
f(\theta) = \frac{1}{k} \sum_{l=0}^{\infty} (2l+1) e^{i\delta_l}\sin(\delta_l) P_l(\cos\theta)
\]
\end{resultbox}

\subsection{חתכי פעולה וריאקציות}

כדי לקבל את חתך הפעולה הכולל, מבצעים אינטגרציה על כל הזוויות ($\int |f(\theta)|^2 d\Omega$). תכונת האורתוגונליות של פולינומי לז'נדר מפשטת את הסכום.

\begin{mainbox}{חתך פעולה לפיזור אלסטי}
\[
\sigma_{el} = \frac{\pi}{k^2} \sum_{l=0}^{\infty} (2l+1) |1 - S_l|^2 = \frac{4\pi}{k^2} \sum_{l=0}^{\infty} (2l+1) \sin^2(\delta_l)
\]
\end{mainbox}

\begin{important}{גבול אנרגיות נמוכות}
כאשר האנרגיה נמוכה ($k \to 0$), היסט הפאזה דועך חזק עם $l$:
\[
\delta_l \sim k^{2l+1}
\]
לכן, באנרגיות נמוכות, הפיזור נשלט כמעט לחלוטין על ידי הגל החלקי $l=0$ (\LR{s-wave}), וניתן להזניח גלים גבוהים יותר.
\end{important}

\subsubsection{ריאקציות (\LR{Inelastic Scattering})}
במקרה שהחלקיק נבלע או משנה מצב (למשל, דיוטרון מתפרק), יש אובדן שטף בערוץ האלסטי. במקרה זה $|S_l| < 1$, ומתקבל חתך פעולה לריאקציה:
\[
\sigma_{reaction}^{(l)} = \frac{\pi}{k^2} (2l+1) (1 - |S_l|^2)
\]
חתך הפעולה המקסימלי לפיזור גדול פי 4 מחתך הפעולה המקסימלי לריאקציה (נובע מכך ש-$S_l$ יכול להיות $-1$ בפיזור, אך מוגבל ל-$0$ בריאקציה מלאה).




\section{תורת הפיזור: המשפט האופטי וגלים חלקיים}

בשיעור זה, שהוא השיעור המסכם של הקורס, נעסוק בשתי תובנות מרכזיות בתורת הפיזור: המשפט האופטי (\LR{The Optical Theorem}) והשיטה של פירוק לגלים חלקיים (\LR{Partial Wave Expansion}). כלים אלו מאפשרים לנו לקשר בין גדלים תיאורטיים (כמו אמפליטודת הפיזור) לגדלים נסיוניים (כמו חתך הפעולה הכולל).

\subsection{המשפט האופטי (\LR{The Optical Theorem})}

המשפט האופטי הוא תוצאה מפתיעה וחזקה, שמקורה באופטיקה (לורד ריילי, סוף המאה ה-19), והיא נובעת משימור הסתברות. המשפט מקשר בין החלק המדומה של אמפליטודת הפיזור הקדמי ($f(\theta=0)$) לבין חתך הפעולה הכולל ($\sigma_{tot}$).

\begin{theorembox}{המשפט האופטי}
עבור גל שנכנס עם וקטור גל $\vec{k}$, חתך הפעולה הכולל נתון על ידי:
\[
\sigma_{total} = \frac{4\pi}{k} \text{Im}[f(0)]
\]
כאשר $f(0)$ היא אמפליטודת הפיזור בזווית $\theta=0$ (פיזור קדימה).
\end{theorembox}

\begin{notebox}{אינטואיציה פיזיקלית}
המשפט נובע מכך שפיזור קדימה מתערבב (interferes) עם הגל המקורי. הירידה בעוצמת הגל המקורי (שמתבטאת בחתך הפעולה הכולל, כלומר כמה חלקיקים הוצאו מהאלומה) חייבת להיות קשורה לאופן שבו גל הפיזור מתאבך עם הגל הפוגע בכיוון הקדמי.
\end{notebox}

\subsubsection{הוכחת המשפט האופטי (מתווה)}

כדי להוכיח את המשפט, נשתמש בקשר בין אמפליטודת הפיזור למטריצת המעבר $T$.

\begin{definitionbox}{אופרטור המעבר T}
אנו מגדירים את הקשר בין הגל החופשי $|k\rangle$ לבין המצב הפיזיקלי המלא $|\psi_{in}\rangle$ (שכולל את הפיזור) באמצעות האופרטור $T$:
\[
T|k\rangle = V|\psi_{in}\rangle
\]
אמפליטודת הפיזור קשורה לאלמנט המטריצה של $T$:
\[
f(\theta=0) = -\frac{m}{2\pi \hbar^2} \langle k | T | k \rangle
\]
\end{definitionbox}

\begin{mainbox}{שלבי ההוכחה}
אנו רוצים לחשב את $\text{Im}\langle k|T|k\rangle$. נשתמש במשוואת ליפמן-שווינגר (\LR{Lippmann-Schwinger}) ובקשרים הבאים:

1. \textbf{הצבת משוואת ליפמן-שווינגר:}
   \[
   |\psi_{in}\rangle = |k\rangle + \frac{1}{E - H_0 + i\epsilon} V |\psi_{in}\rangle
   \]
   נכפיל ב-$\langle k|V$ משמאל (תוך שימוש ב-$T|k\rangle=V|\psi_{in}\rangle$) ונקבל ביטוי עבור $\langle k|T|k\rangle$.

2. \textbf{שימוש בזהות מתמטית:}
   עבור האינטגרל סביב הקוטב (pole), נשתמש בזהות:
   \[
   \frac{1}{x + i\epsilon} = \text{PV}\left(\frac{1}{x}\right) - i\pi \delta(x)
   \]
   החלק של ה-\LR{Principal Value} (\LR{PV}) הוא ממשי (עבור אופרטורים הרמיטיים $V, H_0$) ולכן נופל כאשר לוקחים את החלק המדומה.

3. \textbf{החלק המדומה:}
   החלק המדומה מגיע מהאיבר של הדלתא ($i\pi \delta(...)$). על ידי הכנסת בסיס שלם של מצבי תנע $|\vec{k}'\rangle$, נקבל אינטגרל על כל כיווני הפיזור האפשריים.

4. \textbf{סיכום:}
   האינטגרציה על כל הזוויות נותנת בדיוק את הגדרת חתך הפעולה הכולל $\sigma_{tot}$, וכך מתקבל הקשר:
   \[
   \text{Im} f(0) = \frac{k}{4\pi} \sigma_{tot}
   \]
\end{mainbox}

\subsection{פירוק לגלים חלקיים (\LR{Partial Wave Expansion})}

כאשר הפוטנציאל הוא מרכזי ($V(r)$) וקצר טווח, התנע הזוויתי $L$ נשמר. זה מאפשר לנו לפרק את הבעיה לסכום של בעיות חד-ממדיות עבור כל $L$.

\begin{importantbox}{תנאים לשימוש}
השיטה יעילה כאשר:
1. הפוטנציאל מרכזי (סימטריה כדורית).
2. הפוטנציאל קצר טווח (דועך מהר יותר מ-$1/r$, לא מתאים לפוטנציאל קולומבי טהור ללא התאמות).
\end{importantbox}

\subsubsection{פירוק הגל המישורי}
הגל המישורי הנכנס $e^{i\vec{k}\cdot\vec{r}}$ אינו מצב עצמי של התנע הזוויתי, אך ניתן לפרק אותו לסכום של גלים כדוריים (פונקציות בסל ספיריות $j_l$ ופולינומי לז'נדר $P_l$):

\begin{formulabox}{פירוק גל מישורי (\LR{Rayleigh Expansion})}
\[
e^{i\vec{k}\cdot\vec{r}} = \sum_{l=0}^{\infty} i^l (2l+1) j_l(kr) P_l(\cos\theta)
\]
\end{formulabox}

במרחקים גדולים ($kr \gg 1$), הגל המישורי מתנהג כסופרפוזיציה של גלים כדוריים נכנסים ויוצאים.

\subsubsection{הסחת הפאזה (Phase Shift)}
בנוכחות פוטנציאל, החלקיק המפוזר חייב לשמור על התנע הזוויתי (עבור פוטנציאל מרכזי). בפיזור אלסטי, האמפליטודה לא משתנה, אלא רק הפאזה של הגל היוצא.

\begin{definitionbox}{הסחת הפאזה $\delta_l$ (\LR{Phase Shift})}
אנו מגדירים את הסטייה של הגל היוצא ביחס למקרה החופשי באמצעות זווית הפאזה $\delta_l$.
את השינוי הזה ניתן לתאר באמצעות \textbf{מטריצת הפיזור} (\LR{S-matrix}) עבור גל חלקי $l$:
\[
S_l = e^{2i\delta_l}
\]
עבור פיזור אלסטי $|S_l|=1$, ולכן $\delta_l$ הוא מספר ממשי.
\end{definitionbox}

\subsubsection{חתך הפעולה במונחי גלים חלקיים}

על ידי הצבת הפירוק של פונקציית הגל וחישוב שטף ההסתברות היוצא, ניתן לבטא את אמפליטודת הפיזור ואת חתך הפעולה באמצעות הסחות הפאזה.

\begin{resultbox}{אמפליטודת הפיזור $f(\theta)$}
\[
f(\theta) = \frac{1}{2ik} \sum_{l=0}^{\infty} (2l+1) (S_l - 1) P_l(\cos\theta)
\]
או באמצעות הסחת הפאזה:
\[
f(\theta) = \frac{1}{k} \sum_{l=0}^{\infty} (2l+1) e^{i\delta_l} \sin(\delta_l) P_l(\cos\theta)
\]
\end{resultbox}

מכאן נגזר חתך הפעולה האלסטי הכולל:

\begin{formulabox}{חתך פעולה אלסטי}
\[
\sigma_{el} = \int |f(\theta)|^2 d\Omega = \frac{4\pi}{k^2} \sum_{l=0}^{\infty} (2l+1) \sin^2 \delta_l
\]
\end{formulabox}

\begin{summarybox}{נקודות חשובות}
\begin{itemize}
    \item \textbf{הגבול האוניטרי:} חתך הפעולה המקסימלי עבור גל חלקי $l$ מתקבל כאשר $\delta_l = \frac{\pi}{2}$ (תהודה), והוא שווה ל-$\frac{4\pi}{k^2}(2l+1)$.
    \item \textbf{אנרגיות נמוכות:} בגבול $k \to 0$, הסחת הפאזה מתנהגת כמו $\delta_l \propto k^{2l+1}$. לכן, באנרגיות נמוכות מאוד, רק האיבר $l=0$ (גל-S) דומיננטי, והפיזור הוא איזוטרופי.
    \item \textbf{ריאקציות (פיזור לא אלסטי):} אם יש בליעה (\LR{absorption}), אז $|S_l| < 1$, וניתן להגדיר חתך פעולה לריאקציה ($\sigma_{reaction}$) בנפרד מחתך הפעולה האלסטי.
\end{itemize}
\end{summarybox}

